\newpage
\subsection{Proyección a ambiente productivo}
\label{subsec:produccion}

El prototipo desarrollado está diseñado arquitectónicamente para escalar
hacia un ambiente productivo sin cambios estructurales. Esta sección
presenta un análisis comparativo de las opciones de infraestructura
disponibles, una justificación de las instancias de base de datos
seleccionadas y una estimación de los costos mensuales de infraestructura
en la nube para un escenario productivo con más de 100 dispositivos
móviles activos, más de 200 usuarios en la plataforma web y más de
50,000 imágenes sincronizadas mensualmente.

\subsubsection*{Análisis comparativo: serverless vs contenedores}

Una decisión arquitectónica fundamental para el ambiente productivo es
la elección entre un modelo serverless basado en AWS Lambda y un modelo
basado en contenedores con Docker y Kubernetes (Amazon EKS). La
Tabla~\ref{tab:serverless_vs_eks} resume las diferencias principales
entre ambos enfoques.

\begin{longtable}{|>{\raggedright\arraybackslash}m{3.5cm}|>
{\raggedright\arraybackslash}m{5cm}|>{\raggedright\arraybackslash}m{5cm}|}
\caption{Comparativa entre arquitectura serverless y contenedores para ambiente productivo.}
\label{tab:serverless_vs_eks} \\
\hline
\textbf{Criterio} & \textbf{AWS Lambda (serverless)} & \textbf{Docker + Kubernetes (EKS)} \\
\hline
\endfirsthead
\multicolumn{3}{c}{{\bfseries Tabla \thetable\ -- continuación}} \\
\hline
\textbf{Criterio} & \textbf{AWS Lambda (serverless)} & \textbf{Docker + Kubernetes (EKS)} \\
\hline
\endhead

Administración de infraestructura & Nula, AWS gestiona todo automáticamente & Alta, requiere configuración y mantenimiento del cluster \\
\hline
Escalado & Automático e instantáneo por invocación & Automático pero requiere configuración de HPA y nodos \\
\hline
Cold starts & 2--5 segundos con Micronaut en el primer arranque & Sin cold starts, los contenedores están siempre activos \\
\hline
Costo base mensual & \$0.00 si no hay tráfico & ~\$150 USD solo por el cluster EKS, independiente del tráfico \\
\hline
Load balancer & No necesario, API Gateway lo reemplaza & Necesario (ALB), costo adicional de ~\$20 USD/mes \\
\hline
Límite de ejecución & 15 minutos máximo por función & Sin límite de tiempo de ejecución \\
\hline
Ideal para & Tráfico intermitente o impredecible & Tráfico constante y predecible de alto volumen \\
\hline
Complejidad operativa & Baja & Alta \\
\hline
\end{longtable}

Para el escenario productivo del sistema de detección de insectos
descortezadores, con 100+ dispositivos sincronizando datos de forma
\textbf{ocasional} cuando existe conectividad a internet, el modelo
serverless con AWS Lambda continúa siendo la opción más adecuada. El
tráfico de sincronización es por naturaleza intermitente, que es
exactamente el escenario donde Lambda ofrece mayor ventaja económica
y operativa sobre Kubernetes.

Sin embargo, si el sistema llegara a escalar a \textbf{miles de
dispositivos} con sincronizaciones frecuentes y continuas, o si se
requiriera integrar procesos de larga duración como el análisis masivo
de imágenes en batch, la migración hacia una arquitectura basada en
contenedores con Amazon EKS, un Application Load Balancer (ALB) y
Auto Scaling Groups sería la evolución natural. Esta transición no
requeriría rediseñar la arquitectura base, ya que Micronaut está
diseñado para ejecutarse tanto en Lambda como en contenedores Docker
sin cambios significativos en el código.

\subsubsection*{Análisis comparativo de instancias de base de datos}

La selección de la instancia de Amazon RDS para el ambiente productivo
requiere un balance entre rendimiento, disponibilidad y costo. La
Tabla~\ref{tab:rds_comparativa} presenta un análisis de las principales
opciones disponibles para PostgreSQL en la región \texttt{us-east-1}.

\begin{longtable}{|>{\raggedright\arraybackslash}m{2.5cm}|>
{\centering\arraybackslash}m{1.5cm}|>{\centering\arraybackslash}m{1.8cm}|>
{\centering\arraybackslash}m{2cm}|>{\centering\arraybackslash}m{2cm}|>
{\raggedright\arraybackslash}m{3.5cm}|}
\caption{Comparativa de instancias Amazon RDS PostgreSQL para ambiente productivo (us-east-1, 2026).}
\label{tab:rds_comparativa} \\
\hline
\textbf{Instancia} & \textbf{vCPUs} & \textbf{RAM (GB)} & \textbf{Single-AZ (USD/mes)} & \textbf{Multi-AZ (USD/mes)} & \textbf{Caso de uso} \\
\hline
\endfirsthead
\multicolumn{6}{c}{{\bfseries Tabla \thetable\ -- continuación}} \\
\hline
\textbf{Instancia} & \textbf{vCPUs} & \textbf{RAM (GB)} & \textbf{Single-AZ (USD/mes)} & \textbf{Multi-AZ (USD/mes)} & \textbf{Caso de uso} \\
\hline
\endhead

db.t3.micro & 2 & 1 & \$15.33 & \$30.66 & Desarrollo y pruebas, no recomendado para producción \\
\hline
db.t3.small & 2 & 2 & \$28.49 & \$56.98 & Producción ligera, hasta 50 usuarios concurrentes \\
\hline
\textbf{db.t3.medium} & \textbf{2} & \textbf{4} & \textbf{\$52.56} & \textbf{\$105.12} & \textbf{Producción estándar, 100--200 usuarios, recomendado} \\
\hline
db.t3.large & 2 & 8 & \$105.12 & \$210.24 & Alta concurrencia, 500+ usuarios \\
\hline
db.m5.large & 2 & 8 & \$138.00 & \$276.00 & Producción intensiva, cargas sostenidas sin bursting \\
\hline
\end{longtable}

La instancia \textbf{db.t3.medium} fue seleccionada para el ambiente
productivo por las siguientes razones:

\begin{itemize}
    \item \textbf{Memoria suficiente:} Con 4 GB de RAM puede mantener
    en caché las consultas más frecuentes del sistema, como el historial
    de detecciones y los registros de árboles, reduciendo la latencia
    de respuesta.
    \item \textbf{CPU con bursting:} La familia T3 opera en modo
    \textit{unlimited}, permitiendo ráfagas de CPU cuando la carga
    aumenta durante las sincronizaciones masivas, sin requerir una
    instancia de mayor tamaño de forma permanente.
    \item \textbf{Costo equilibrado:} Representa el punto óptimo entre
    rendimiento y costo para el volumen estimado de 100+ dispositivos
    y 200+ usuarios, con reservación anual para reducir el costo
    efectivo mensual \cite{aws_rds_pricing}.
    \item \textbf{Escalabilidad hacia Aurora:} Si el sistema escalara
    a miles de dispositivos, la migración hacia \textbf{Amazon Aurora
    PostgreSQL} no requeriría cambios en el código del backend, ya
    que Aurora es compatible con el protocolo PostgreSQL.
\end{itemize}

\subsubsection*{Estimación de costos mensuales}

La siguiente estimación fue generada mediante la herramienta oficial
AWS Pricing Calculator \cite{aws_calculator} para la región
\texttt{us-east-1} (US East - N. Virginia), vigente a mayo de 2026.
El tipo de cambio utilizado para la conversión es de \$17.50 MXN por
dólar estadounidense (USD).

\subsubsection*{Supuestos del escenario productivo}

\begin{itemize}
    \item \textbf{Dispositivos móviles activos:} 100+ dispositivos
    sincronizando datos de forma periódica.
    \item \textbf{Usuarios plataforma web:} 200+ usuarios con sesiones
    activas.
    \item \textbf{Imágenes sincronizadas:} 50,000+ imágenes/mes con
    un tamaño promedio de 2 MB por imagen, equivalente a 100 GB de
    almacenamiento mensual nuevo.
    \item \textbf{Invocaciones Lambda:} 5 millones de invocaciones/mes
    con arquitectura ARM64 (Graviton2), 500ms de duración promedio y
    512 MB de memoria asignada.
    \item \textbf{Base de datos:} Instancia \texttt{db.t3.medium}
    PostgreSQL con Reserved Instance de 1 año y pago parcial anticipado,
    100 GB de almacenamiento SSD gp3.
\end{itemize}

\begin{longtable}{|>{\raggedright\arraybackslash}m{4cm}|>
{\raggedright\arraybackslash}m{4.5cm}|>{\centering\arraybackslash}m{2cm}|>
{\centering\arraybackslash}m{2cm}|>{\centering\arraybackslash}m{2cm}|}
\caption{Estimación de costos en ambiente productivo generada con AWS Pricing Calculator (us-east-1, 2026).}
\label{tab:costos_produccion} \\
\hline
\textbf{Servicio} & \textbf{Configuración} & \textbf{Upfront (USD)} & \textbf{Mensual (USD)} & \textbf{Mensual (MXN)} \\
\hline
\endfirsthead
\multicolumn{5}{c}{{\bfseries Tabla \thetable\ -- continuación}} \\
\hline
\textbf{Servicio} & \textbf{Configuración} & \textbf{Upfront (USD)} & \textbf{Mensual (USD)} & \textbf{Mensual (MXN)} \\
\hline
\endhead

AWS Lambda & 5M invocaciones, ARM64, 512 MB, 500ms & \$0.00 & \$0.80 & \$14.00 \\
\hline
Amazon API Gateway & 5M llamadas REST/mes & \$0.00 & \$17.50 & \$306.25 \\
\hline
Amazon RDS PostgreSQL & db.t3.medium, Reserved 1 año, 100 GB gp3 & \$845.00 & \$63.15 & \$1,105.13 \\
\hline
Amazon S3 & 500 GB almacenamiento estándar & \$0.00 & \$11.83 & \$206.93 \\
\hline
Amazon CloudFront & Pro Plan (1 distribución) & \$0.00 & \$15.00 & \$262.50 \\
\hline
Amazon CloudWatch & Logs y métricas & \$0.00 & \$3.00 & \$52.50 \\
\hline
AWS WAF & 1 Web ACL, reglas básicas & \$0.00 & \$13.00 & \$227.50 \\
\hline
Amazon Cognito & 200+ usuarios activos/mes & \$0.00 & \$10.00 & \$175.00 \\
\hline
\textbf{Total} & & \textbf{\$845.00} & \textbf{\$134.28} & \textbf{\$2,349.90} \\
\hline
\end{longtable}

El costo mensual recurrente en ambiente productivo asciende a
\textbf{\$134.28 USD (\$2,349.90 MXN)}, con un pago único inicial de
\textbf{\$845.00 USD (\$14,787.50 MXN)} correspondiente a la
reservación anual de la instancia Amazon RDS. El costo total estimado
para los primeros 12 meses de operación en producción es de
\textbf{\$2,456.36 USD (\$42,986.30 MXN)}, incluyendo el pago
anticipado de RDS.

La reservación anual de RDS representa un ahorro significativo respecto
al precio bajo demanda: la instancia \texttt{db.t3.medium} en
configuración on-demand costaría aproximadamente \$105 USD/mes
(\$1,260 USD/año), mientras que con Reserved Instance el costo
efectivo se reduce a \$845 + (\$63.15 $\times$ 12) = \$1,602.80
USD/año, un ahorro de alrededor del \textbf{27\%} respecto al precio
on-demand \cite{aws_rds_pricing}.

\subsubsection*{Estrategias de optimización de costos}

Para reducir el costo mensual en producción, se pueden aplicar las
siguientes estrategias dentro del mismo ecosistema AWS:

\begin{itemize}
    \item \textbf{Reserved Instances en RDS:} Ya aplicado en esta
    estimación, con un ahorro del 27\% respecto al precio on-demand.
    \item \textbf{Lambda ARM64 (Graviton2):} Las funciones Lambda
    configuradas con arquitectura ARM64 ofrecen un 20\% de reducción
    en costos de cómputo sin cambios en el código de Micronaut.
    \item \textbf{S3 Intelligent-Tiering:} Para imágenes antiguas
    con acceso poco frecuente, S3 Intelligent-Tiering reduce
    automáticamente el costo de almacenamiento hasta un 40\%.
    \item \textbf{CloudFront caché:} Configurar políticas de caché
    agresivas para los artefactos Angular reduce las solicitudes
    al origen S3, disminuyendo tanto el costo de CloudFront como
    el de S3.
\end{itemize}

La estimación completa puede consultarse en el siguiente enlace
generado por AWS Pricing Calculator:

\begin{center}
\url{https://calculator.aws/#/estimate?id=65c11c7fb1e7ea52aab349e587a999504e736112}
\end{center}
\cite{aws_calculator}